\documentclass[UTF8,a4paper,12pt,oneside]{ctexrep}

\usepackage{geometry}
\usepackage{setspace}
\usepackage{booktabs}
\usepackage{longtable}
\usepackage{tabularx}
\usepackage{array}
\usepackage{enumitem}
\usepackage{xcolor}
\usepackage{hyperref}
\usepackage{fancyhdr}
\usepackage{titlesec}
\usepackage{float}
\usepackage{graphicx}

\geometry{left=2.8cm,right=2.8cm,top=2.8cm,bottom=2.8cm}
\onehalfspacing
\setlength{\parindent}{2em}
\setlength{\parskip}{0.25em}
\setlist[itemize]{leftmargin=2.5em}
\setlist[enumerate]{leftmargin=2.7em}

\hypersetup{
  colorlinks=true,
  linkcolor=black,
  urlcolor=blue,
  citecolor=black,
  pdfauthor={seccomp-privacy-platform team},
  pdftitle={面向电商隐私分析场景的 SSE + PJC 平台设计与实现}
}

\pagestyle{fancy}
\setlength{\headheight}{15pt}
\fancyhf{}
\fancyhead[C]{面向电商隐私分析场景的 SSE + PJC 平台设计与实现}
\fancyfoot[C]{\thepage}
\renewcommand{\headrulewidth}{0.4pt}

\ctexset{
  chapter = {
    name = {第,章},
    number = \chinese{chapter},
    format = \centering\zihao{3}\heiti,
    beforeskip = 12pt,
    afterskip = 18pt
  },
  section = {
    format = \zihao{4}\heiti,
    beforeskip = 12pt,
    afterskip = 6pt
  },
  subsection = {
    format = \zihao{-4}\heiti,
    beforeskip = 8pt,
    afterskip = 4pt
  }
}

\newcommand{\project}{seccomp-privacy-platform}
\newcommand{\filepath}[1]{\texttt{\detokenize{#1}}}
\newcommand{\placeholder}[1]{\textcolor{gray}{#1}}

\begin{document}

\begin{titlepage}
  \centering
  \vspace*{1.5cm}
  {\zihao{2}\heiti 面向电商隐私分析场景的\\[0.3em]SSE + PJC 平台设计与实现\par}
  \vspace{1.2cm}
  {\zihao{3}\bfseries 项目论文模板（提纲版）\par}
  \vspace{2.2cm}
  \begin{tabular}{>{\bfseries}r l}
    作品名称： & \underline{\hspace{7cm}} \\
    项目仓库： & \project \\
    作者： & \underline{\hspace{7cm}} \\
    指导教师： & \underline{\hspace{7cm}} \\
    学校/学院： & \underline{\hspace{7cm}} \\
    日期： & \today \\
  \end{tabular}
  \vfill
  {\zihao{-4}\placeholder{说明：本模板只提供论文结构、模块分类和写作占位，不预写正文。具体内容由负责论文撰写的同学补充。}\par}
\end{titlepage}

\pagenumbering{Roman}

\chapter*{摘要}
\addcontentsline{toc}{chapter}{摘要}
\placeholder{在此填写摘要。建议包含：研究背景、核心问题、技术路线（database/control-plane + SSE + Google PJC）、系统实现、主要结果、结论。}

\vspace{1em}
\noindent\textbf{关键词：}\placeholder{可搜索加密；PJC；隐私计算；电商分析；审计治理}

\tableofcontents
\clearpage
\pagenumbering{arabic}

\chapter{引言}

\section{研究背景}
\placeholder{说明为什么电商场景需要跨方隐私分析，传统明文交换或普通访问控制为什么不够。}

\section{问题定义}
\placeholder{明确本文要解决的问题：在不暴露原始标识和业务明细的前提下，如何完成候选筛选、交集计算、结果发布和审计治理。}

\section{项目定位}
\placeholder{说明本项目不是完整电商平台，而是面向电商隐私分析场景的 SSE + PJC 平台化实现。}

\section{本文贡献}
\placeholder{列出 3--5 条贡献，例如：主链路实现、平台治理、typed schema、workflow、evidence gate、业务场景适配等。}

\chapter{系统总体设计}

\section{总体技术主线}
\placeholder{建议给出主链路文字说明：metadata/control-plane $\rightarrow$ SSE $\rightarrow$ recovery $\rightarrow$ bridge $\rightarrow$ PJC $\rightarrow$ release。}

\section{模块分类}

这一节建议作为全文的“项目基石”部分，用来建立完整项目结构。后续章节都围绕这些模块展开。

\begin{longtable}{p{0.22\textwidth}p{0.68\textwidth}}
\toprule
\textbf{模块} & \textbf{职责} \\
\midrule
\filepath{sse/} & \placeholder{填写 SSE 模块职责：密文检索、候选导出、策略约束等。} \\
\filepath{services/record_recovery/} & \placeholder{填写记录恢复模块职责：最小必要字段恢复、服务边界、authz、health 等。} \\
\filepath{bridge/} & \placeholder{填写 bridge 模块职责：规范化、token 化、PJC 输入准备、input commitment。} \\
\filepath{a-psi/} & \placeholder{填写 Google PJC 集成职责：交集统计、结果聚合、发布前结果承接。} \\
\filepath{scripts/} & \placeholder{填写 orchestration / workflow / gate / benchmark / archive 等脚本层职责。} \\
\filepath{migrations/metadata/} & \placeholder{填写控制面数据库职责：workflow、policy、permission、audit、privacy budget。} \\
\filepath{schemas/} & \placeholder{填写 typed contract / verifier-facing schema 体系职责。} \\
\filepath{docs/} & \placeholder{填写设计文档、威胁模型、治理文档、runbook 职责。} \\
\bottomrule
\end{longtable}

\section{总体架构图}
\placeholder{此处插入总体架构图。建议包含主链路和控制面/审计/发布治理侧能力。}

\begin{figure}[H]
\centering
\fbox{\parbox[c][6cm][c]{0.85\textwidth}{\centering \placeholder{在此插入总体架构图或流程图}}}
\caption{系统总体架构图（占位）}
\end{figure}

\chapter{关键模块设计与实现}

\section{SSE 模块}
\subsection{设计目标}
\placeholder{写 SSE 为什么存在、解决什么问题。}
\subsection{核心机制}
\placeholder{写 searchable encryption、candidate selection、policy-gated export。}
\subsection{当前边界}
\placeholder{说明 legacy SSE WebSocket 已退休，核心保留的是 candidate selection 能力。}

\section{Record Recovery 模块}
\subsection{设计目标}
\placeholder{写为什么需要受控记录恢复。}
\subsection{实现机制}
\placeholder{写 encrypted record store、candidate-bound recovery、service boundary。}
\subsection{安全意义}
\placeholder{写它如何减少明文扩散。}

\section{Bridge 模块}
\subsection{设计目标}
\placeholder{写为什么要做 join key normalization 和 tokenization。}
\subsection{实现机制}
\placeholder{写 HMAC token、job metadata、input commitment、handoff。}
\subsection{安全意义}
\placeholder{写 bridge 在原始标识与 PJC 输入之间承担的边界作用。}

\section{Google PJC 模块}
\subsection{设计目标}
\placeholder{写 PJC 在项目中的职责。}
\subsection{实现机制}
\placeholder{写 server/client 输入、交集统计、结果生成。}
\subsection{边界说明}
\placeholder{明确说明当前只能表述为 semi-honest/operator-controlled。}

\section{Control Plane 与 Query Workflow}
\subsection{Metadata Sidecar}
\placeholder{写控制面数据库承载的实体和状态。}
\subsection{Query Workflow}
\placeholder{写 request / submit / enqueue / worker / release 的闭环。}
\subsection{Operator Surface}
\placeholder{写 operator dashboard、审批、request workflow、public summary。}

\chapter{安全设计与治理机制}

\section{安全目标}
\placeholder{写本项目重点保护什么，不保护什么。}

\section{威胁模型}
\placeholder{写 semi-honest、query-abuse、malicious participant、operator-side 风险等。}

\section{协议内与协议外边界}
\placeholder{明确区分 protocol-internal 与 protocol-external governance。}

\section{当前已实现的安全增强}
\placeholder{写 input commitment、source attestation、release gate、privacy budget、audit chain、public summary 等。}

\section{当前仍然保留的边界}
\placeholder{写 malicious-secure 未解决、外部 trust-root 依赖、SSE 更强泄露模型未解决等。}

\chapter{电商叙事环境与业务适配}

\section{为什么选择电商场景}
\placeholder{说明订单、归因、支付、履约、客服等场景为什么适合作为隐私计算平台叙事环境。}

\section{电商事实层适配}
\placeholder{介绍事实表、业务身份、persona、ABAC/workbench 等。}

\section{电商叙事与技术主线的关系}
\placeholder{强调电商是业务叙事环境，不是技术主线替代品。}

\chapter{系统验证与实验结果}

\section{验证方法}
\placeholder{写 schema contract、smoke、negative tests、benchmark、deployment evidence、closure gate。}

\section{关键结果}
\placeholder{填写关键数字、关键截图、关键 gate 结果。}

\section{模块级状态汇总}
\placeholder{按模块写当前 repo-side / live-side 状态。}

\section{客观评价}
\placeholder{这一节建议从工程完成度、协议安全边界、平台化程度、答辩价值四个角度评价。}

\chapter{不足与改进方向}

\section{当前不足}
\placeholder{写协议内边界、外部 trust-root 边界、长期运维边界。}

\section{不更换协议前提下的改进方向}
\placeholder{写 strict source attestation、release gate、privacy budget、workflow governance 等。}

\section{必须超出当前框架才能解决的问题}
\placeholder{写 malicious-secure backend、active-secure 计算、更强 SSE 泄露缓解、外部 trust-root。}

\chapter{结论}
\placeholder{总结项目定位、主要实现、平台价值、边界与意义。}

\appendix

\chapter{写作分工建议}

\begin{enumerate}
  \item 一名同学负责“引言、背景、问题定义、相关工作”
  \item 一名同学负责“系统设计、模块实现、业务适配”
  \item 一名同学负责“安全设计、威胁模型、验证结果、结论”
\end{enumerate}

\chapter{插图与表格建议}

\begin{enumerate}
  \item 总体架构图
  \item 主链路流程图
  \item 模块职责表
  \item 安全边界示意图
  \item 模块状态汇总表
  \item 关键实验/验证结果截图
\end{enumerate}

\chapter{可替换字段}

\begin{itemize}
  \item 标题、作者、学校、指导教师
  \item 摘要、关键词
  \item 架构图、流程图、截图
  \item 当前模块状态数字
  \item 评估结论和不足部分
\end{itemize}

\end{document}
